\PilotPage{19}
\begin{equation}
    s^2=-\omega_n^2.
    \tag{13}
\end{equation}
Eq.~(13) has two imaginary roots
\begin{equation}
    s_1=+i\omega_n, \qquad s_2=-i\omega_n.
    \tag{14}
\end{equation}

\subsubsection{General Solution}
Both $s_1=i\omega_n$ and $s_2=-i\omega_n$ yield solutions of the ODE Eq.~(4). These solutions are in the form of complex exponentials $e^{i\omega_n t}$ and $e^{-i\omega_n t}$, respectively. Hence, the general solution of the linear ODE Eq.~(4) is a linear superposition of the two individual solutions, i.e.,
\begin{equation}
    u(t)=C_1e^{i\omega_n t}+C_2e^{-i\omega_n t}
    \tag{15}
\end{equation}
where $C_1$ and $C_2$ are arbitrary constants that need to be determined from the initial conditions.

Recall Euler's identity
\begin{equation}
    e^{i\alpha}=\cos\alpha+i\sin\alpha
    \tag{16}
\end{equation}
and hence
\begin{equation}
    e^{-i\alpha}=\cos\alpha-i\sin\alpha.
    \tag{17}
\end{equation}
Eqs.~(16), (17) can be used to rewrite Eq.~(15) in trigonometric form, i.e.,
\begin{equation}
    u(t)=A\cos\omega_n t+B\sin\omega_n t.
    \tag{18}
\end{equation}
where the constants $A$, $B$ and $C_1$, $C_2$ are interconnected as
\begin{align}
    A&=C_1+C_2, \tag{19}\\
    B&=i(C_1-C_2). \tag{20}
\end{align}
Further processing allows us to write Eq.~(18) as
\begin{equation}
    u(t)=C\sin(\omega_n t+\varphi)
    \tag{21}
\end{equation}
where
\begin{equation}
    C=\sqrt{A^2+B^2}, \qquad \varphi=\tan^{-1}\frac{A}{B}.
    \tag{22}
\end{equation}

\subsubsection{Initial Value Problem (IVP)}
The initial value problem (IVP) consists in finding the general-solution constants in terms of the initial conditions of the vibrating system, i.e., initial displacement $u_0$ and initial velocity $\dot{u}_0$. The purpose is to find the constants $A$, $B$ in terms of known $u_0$ and $\dot{u}_0$. Using Eq.~(18), we write
